\section{Consequences and limits}\label{sec:scope}
The results separate the size of a copositive certificate from the number of adjacent pivots needed to reach one. A finite positive reserve turns each certified seed into robust lower bounds for every improving choice from a bounded region of one fixed-block face; the common separator remains short. Independently, the fixed-block epigraph has finite local configurations and intrinsic graph growth of exact degree equal to the block order. Graphically different seeds, including branched leading blocks and principal-minimal nonmembers, show how the amplification extends beyond the five-cycle. The excluded methods are precise: unrelated global certificates, distant input-dependent restarts, nonadjacent jumps, and non-descending steps. Bounded-region starts with different leading perfect blocks are also outside the theorem; the same strict normal reserve cannot favor two distinct perfect vertices. Intrinsic graph distance does not control shortcuts outside the face. None of these results decides eventual termination under a suitable pivot rule or gives a general polynomial-time recognition algorithm. Exact-arithmetic programs and certificate data accompany the paper as supplementary material.
